% config/thesis-config.tex
% -----------------------------------------------------------------------------
% USER CONFIGURATION
% -----------------------------------------------------------------------------
% Change only this value to switch labels and selected content folders.
% The template loads files from content/<language>, chapters/<language>,
% and appendices/<language>. Figures are looked up in figures/english/ first
% when the language is english, then in figures/.
% Valid values: indonesian, english
% \providecommand lets the English drivers (main-en.tex, main-en-cetak.tex)
% set the language before this file is read, without editing it here.
\providecommand{\ThesisLanguage}{indonesian}

% Change this value to choose the document type.
% Valid values: skripsi, tesis
% English labels: skripsi -> Undergraduate Thesis, tesis -> Thesis
\newcommand{\ThesisDocumentType}{skripsi}

% Visibility of the writing-guide boxes (\TemplateTip, \TemplateTodo) and the
% draft markers (\Placeholder, \FigurePlaceholder).
%   auto -> shown in main.tex and its print/English drivers, hidden in
%           proposal.tex, skripsi-draft.tex, and semhas.tex
%   on   -> always shown
%   off  -> always hidden (use this for the final submitted manuscript)
% When hidden, \FigurePlaceholder collapses to nothing but the surrounding
% figure keeps its caption, number, and label.
\newcommand{\ThesisShowHelpers}{auto}

% Indonesian and English titles. Both are kept because Indonesian university
% documents commonly need both title forms even when the manuscript language changes.
\newcommand{\ThesisTitleID}{Judul Skripsi Bahasa Indonesia Ditulis dengan Huruf Kapital di Setiap Awal Kata}
\newcommand{\ThesisTitleEN}{English Thesis Title Written in Title Case}

\newcommand{\StudentName}{Nama Lengkap Mahasiswa}
\newcommand{\StudentID}{Nomor Pokok Mahasiswa}
\newcommand{\ProgramNameID}{Informatika}
\newcommand{\ProgramNameEN}{Informatics}

\newcommand{\AdvisorOne}{Nama Dosen Pembimbing I, Gelar}
\newcommand{\AdvisorOneID}{NIP/NPT. \makebox[4cm]{\dotfill}}
\newcommand{\AdvisorTwo}{Nama Dosen Pembimbing II, Gelar}
\newcommand{\AdvisorTwoID}{NIP/NPT. \makebox[4cm]{\dotfill}}
\newcommand{\ExaminerOne}{Nama Ketua Penguji, Gelar}
\newcommand{\ExaminerOneID}{NIP/NPT. \makebox[4cm]{\dotfill}}
\newcommand{\ExaminerTwo}{Nama Penguji I, Gelar}
\newcommand{\ExaminerTwoID}{NIP/NPT. \makebox[4cm]{\dotfill}}

% Official metadata retained from the supplied template. Verify with the latest
% faculty/program guide before official submission.
\newcommand{\DeanName}{Nama Dekan Fakultas, Gelar}
\newcommand{\DeanID}{NIP./NPT. \makebox[4cm]{\dotfill}}
\newcommand{\HeadOfProgramName}{Nama Koordinator Program Studi, Gelar}
\newcommand{\HeadOfProgramID}{NIP./NPT. \makebox[4cm]{\dotfill}}

\newcommand{\DefenseDateID}{Tanggal Bulan Tahun}
\newcommand{\DefenseDateEN}{Month Day, Year}
% Signing date of PERNYATAAN ORISINALITAS and KATA PENGANTAR. It is usually a
% few days after the defense, once revisions are done.
\newcommand{\SignDateID}{Tanggal Bulan Tahun}
\newcommand{\SignDateEN}{Month Day, Year}

% Link encoded as the QR code beside the signature on PERNYATAAN ORISINALITAS,
% as in the faculty sample book. Point it at the supporting files the study
% program asks for, for example a shared folder. Leave the braces empty to
% omit the QR code.
\newcommand{\OriginalityQRLink}{https://example.com/ganti-dengan-tautan-berkas-pendukung}

\newcommand{\SubmissionYear}{2026}
\newcommand{\KeywordsID}{kata kunci pertama, kata kunci kedua, kata kunci ketiga}
\newcommand{\KeywordsEN}{first keyword, second keyword, third keyword}

% Orange outer cover printed before the white inner cover.
%   auto -> only in the final manuscript (main.tex and its print/English drivers)
%   on   -> every manuscript
%   off  -> never (use for files that only need the white cover)
% \providecommand lets a driver override it, for example a *-nocover driver.
\providecommand{\ThesisColorCover}{auto}
